\begin{table}[!htbp]
\centering\small
\caption{Method names, computational identifiers, and predictor dimensions.}
\label{tab:crosswalk}
\begin{tabularx}{\textwidth}{@{}l>{\raggedright\arraybackslash}Xrr@{}}
\toprule
Identifier & Method name and composition & EEG & Total\\
\midrule
M0 & Demographics & 0 & 2\\
M1 & Demographics and task behavior & 0 & 5\\
M2 & Non-EEG baseline: M1 plus object gaze & 0 & 7\\
adaptive M3 & Primary EEG-augmented procedure; selects E0, E1, or E2 & varies & varies\\
E0 & Regional spectral features, added to M2 & 12 & 19\\
E1 & Expanded spectral features: E0 plus topography, variability, spectral entropy & 35 & 42\\
E2 & Expanded spectral and covariance features: E1 plus covariance components & 47 & 54\\
E3 & Phase-lag connectivity, added to M2 & 9 & 16\\
E4 & Permutation entropy and Hjorth features, added to M2 & 9 & 16\\
E5 & Spectral-state features, added to M2 & 9 & 16\\
\bottomrule
\end{tabularx}
\par\smallskip\footnotesize EEG counts are final dimensions after any training-only reduction. E0--E2 form the preregistered library; E3--E5 are separate post-hoc families. Raw dimensions and transformations are specified below.
\end{table}
